\documentclass[10pt,a4paper]{amsart}
\usepackage{amsmath,amssymb,amsthm,graphicx}
\usepackage{verbatim}
\newenvironment{reference}{\comment}{\endcomment}
%\newenvironment{reference}{}{}
\newenvironment{slogan}{\comment}{\endcomment}
\newenvironment{history}{\comment}{\endcomment}
% For commutative diagrams we use Xy-pic
\usepackage[all]{xy}

% We use 2cell for 2-commutative diagrams.
\xyoption{2cell}
% Theorem environments.
%
\theoremstyle{plain}
\newtheorem{theorem}[subsection]{Theorem}
\newtheorem{proposition}[subsection]{Proposition}
\newtheorem{lemma}[subsection]{Lemma}

\theoremstyle{definition}
\newtheorem{definition}[subsection]{Definition}
\newtheorem{example}[subsection]{Example}
\newtheorem{exercise}[subsection]{Exercise}
\newtheorem{situation}[subsection]{Situation}

\theoremstyle{remark}
\newtheorem{remark}[subsection]{Remark}
\newtheorem{remarks}[subsection]{Remarks}

\numberwithin{equation}{subsection}

% Macros
%
\def\lim{\mathop{\mathrm{lim}}\nolimits}
\def\colim{\mathop{\mathrm{colim}}\nolimits}
\def\Spec{\mathop{\mathrm{Spec}}}
\def\Hom{\mathop{\mathrm{Hom}}\nolimits}
\def\Ext{\mathop{\mathrm{Ext}}\nolimits}
\def\SheafHom{\mathop{\mathcal{H}\!\mathit{om}}\nolimits}
\def\SheafExt{\mathop{\mathcal{E}\!\mathit{xt}}\nolimits}
\def\Sch{\mathit{Sch}}
\def\Mor{\mathop{\mathrm{Mor}}\nolimits}
\def\Ob{\mathop{\mathrm{Ob}}\nolimits}
\def\Sh{\mathop{\mathit{Sh}}\nolimits}
\def\NL{\mathop{N\!L}\nolimits}
\def\CH{\mathop{\mathrm{CH}}\nolimits}
\def\proetale{{pro\text{-}\acute{e}tale}}
\def\etale{{\acute{e}tale}}
\def\QCoh{\mathit{QCoh}}
\def\Ker{\mathop{\mathrm{Ker}}}
\def\Im{\mathop{\mathrm{Im}}}
\def\Coker{\mathop{\mathrm{Coker}}}
\def\Coim{\mathop{\mathrm{Coim}}}

% Boxtimes
%
\DeclareMathSymbol{\boxtimes}{\mathbin}{AMSa}{"02}

%
% Macros for moduli stacks/spaces
%
\def\QCohstack{\mathcal{QC}\!\mathit{oh}}
\def\Cohstack{\mathcal{C}\!\mathit{oh}}
\def\Spacesstack{\mathcal{S}\!\mathit{paces}}
\def\Quotfunctor{\mathrm{Quot}}
\def\Hilbfunctor{\mathrm{Hilb}}
\def\Curvesstack{\mathcal{C}\!\mathit{urves}}
\def\Polarizedstack{\mathcal{P}\!\mathit{olarized}}
\def\Complexesstack{\mathcal{C}\!\mathit{omplexes}}
% \Pic is the operator that assigns to X its picard group, usage \Pic(X)
% \Picardstack_{X/B} denotes the Picard stack of X over B
% \Picardfunctor_{X/B} denotes the Picard functor of X over B
\def\Pic{\mathop{\mathrm{Pic}}\nolimits}
\def\Picardstack{\mathcal{P}\!\mathit{ic}}
\def\Picardfunctor{\mathrm{Pic}}
\def\Deformationcategory{\mathcal{D}\!\mathit{ef}}
\usepackage{cleveref}
\crefname{originallemma}{original lemma}{original lemmas}
\Crefname{originallemma}{Original lemma}{Original lemmas}
\makeatletter
\newlabel{lemma-cover-upstairs}{{lemma-cover-upstairs}{source}{}{sourcepointer-1}{}}
\newlabel{lemma-cover-upstairs@cref}{{[originallemma][0][]lemma-cover-upstairs}{[1][0][]source}{}{}{}}
\makeatother
\crefname{originallemma}{original lemma}{original lemmas}
\Crefname{originallemma}{Original lemma}{Original lemmas}
\makeatletter
\newlabel{lemma-factor-factor}{{lemma-factor-factor}{source}{}{sourcepointer-2}{}}
\newlabel{lemma-factor-factor@cref}{{[originallemma][0][]lemma-factor-factor}{[1][0][]source}{}{}{}}
\makeatother
\crefname{originalreference}{original reference}{original references}
\Crefname{originalreference}{Original reference}{Original references}
\makeatletter
\newlabel{morphisms-lemma-flat-base-change-scheme-theoretic-image}{{morphisms-lemma-flat-base-change-scheme-theoretic-image}{source}{}{sourcepointer-3}{}}
\newlabel{morphisms-lemma-flat-base-change-scheme-theoretic-image@cref}{{[originalreference][0][]morphisms-lemma-flat-base-change-scheme-theoretic-image}{[1][0][]source}{}{}{}}
\makeatother
\crefname{originalreference}{original reference}{original references}
\Crefname{originalreference}{Original reference}{Original references}
\makeatletter
\newlabel{morphisms-section-scheme-theoretic-image}{{morphisms-section-scheme-theoretic-image}{source}{}{sourcepointer-4}{}}
\newlabel{morphisms-section-scheme-theoretic-image@cref}{{[originalreference][0][]morphisms-section-scheme-theoretic-image}{[1][0][]source}{}{}{}}
\makeatother
\crefname{originalreference}{original reference}{original references}
\Crefname{originalreference}{Original reference}{Original references}
\makeatletter
\newlabel{spaces-morphisms-lemma-flat-base-change-scheme-theoretic-image}{{spaces-morphisms-lemma-flat-base-change-scheme-theoretic-image}{source}{}{sourcepointer-5}{}}
\newlabel{spaces-morphisms-lemma-flat-base-change-scheme-theoretic-image@cref}{{[originalreference][0][]spaces-morphisms-lemma-flat-base-change-scheme-theoretic-image}{[1][0][]source}{}{}{}}
\makeatother
\crefname{originalreference}{original reference}{original references}
\Crefname{originalreference}{Original reference}{Original references}
\makeatletter
\newlabel{spaces-morphisms-section-scheme-theoretic-image}{{spaces-morphisms-section-scheme-theoretic-image}{source}{}{sourcepointer-6}{}}
\newlabel{spaces-morphisms-section-scheme-theoretic-image@cref}{{[originalreference][0][]spaces-morphisms-section-scheme-theoretic-image}{[1][0][]source}{}{}{}}
\makeatother
\crefname{originalreference}{original reference}{original references}
\Crefname{originalreference}{Original reference}{Original references}
\makeatletter
\newlabel{stacks-properties-lemma-substacks-presentation}{{stacks-properties-lemma-substacks-presentation}{source}{}{sourcepointer-7}{}}
\newlabel{stacks-properties-lemma-substacks-presentation@cref}{{[originalreference][0][]stacks-properties-lemma-substacks-presentation}{[1][0][]source}{}{}{}}
\makeatother
\title{Formation of the scheme theoretic image of a quasi-compact morphism commutes with flat base change}
\author{The Stacks Project Authors}
\date{stacks-project:stacks-morphisms.tex:local-snapshot-20261012; GNU FDL 1.2 or later}
\begin{document}
\maketitle
\noindent\textit{Editorial note.} All delivered mathematical text is the author's own, copied from the pinned Stacks Project chapter source sources/stacks-algebraic-stacks-morphisms/pinned-original-source.tex; no statement or proof step was rewritten, reordered or removed.\par\medskip
\noindent\textit{Editorial note.} Preamble provenance: the Stacks Project keeps one single shared preamble for every chapter; this card takes it directly from the staged sources/stacks-algebraic-stacks-morphisms/preamble.tex as three verbatim spans (lines 8-12, 14-18 and 154-217), so no derived file was written and nothing was added or redefined. The wrapper class and the base packages are supplied by the assembler.\par\medskip
\noindent\textit{Editorial note.} Bibliography provenance: neither the counted statement nor its proof contains a citation (this lane rescans the whole span before writing the card), so the card carries no bibliography and no entry was reconstructed.\par\medskip
\noindent\textit{Editorial note.} Selection provenance: this card was opened under the file-based control contract (CONTROL.md), with the mechanical pre-card gate run on its statement and proof interval before the card was written.\par\medskip
\noindent\textit{Allowed standard prerequisites (editorial).}\begin{itemize}
\item The scheme theoretic image of a morphism of algebraic stacks, and its description through closed substacks and through smooth presentations. The author's section on scheme theoretic images, in which the counted statement occurs. Fixes the four schemes and stacks whose equality is to be proved.
\item The correspondence between closed substacks of a stack and invariant closed subspaces of a smooth presentation, and the invariance of the scheme theoretic image under the relation. Invoked inside the counted proof (source lines 8921-8932) as a cross-chapter pointer. Reduces the comparison of closed substacks to the invariance of a scheme theoretic image.
\item The behaviour of the scheme theoretic image under flat base change in the case of schemes, and the invariance of the image under a further surjective smooth cover. Both are used inside the counted proof (source lines 8947-8948 and 8954-8956, 8980-8983), the first with a pointer to the author lemma lemma-cover-upstairs counted separately in this lane. Allows the image to be computed after replacing the morphism by a surjective smooth covering morphism of schemes, and to be compared after flat base change.
\end{itemize}\medskip
\ifdefined\hypertarget\hypertarget{sourcepointer-1}{}\fi
\noindent\textit{Source reference.} Original lemma lemma-cover-upstairs: sources/stacks-algebraic-stacks-morphisms/pinned-original-source.tex, source lines 8790--8790. Navigation pointer only: the counted text refers to the Stacks item lemma-cover-upstairs, which belongs to another chapter of the project and is not part of this card. The pointer carries the target's own label as its number because this chapter's staged snapshot does not contain the other chapter's tag table; its content is not reproduced here.\par\medskip
\ifdefined\hypertarget\hypertarget{sourcepointer-2}{}\fi
\noindent\textit{Source reference.} Original lemma lemma-factor-factor: sources/stacks-algebraic-stacks-morphisms/pinned-original-source.tex, source lines 8916--8916. Navigation pointer only: the counted text refers to the Stacks item lemma-factor-factor, which belongs to another chapter of the project and is not part of this card. The pointer carries the target's own label as its number because this chapter's staged snapshot does not contain the other chapter's tag table; its content is not reproduced here.\par\medskip
\ifdefined\hypertarget\hypertarget{sourcepointer-3}{}\fi
\noindent\textit{Source reference.} Original reference morphisms-lemma-flat-base-change-scheme-theoretic-image: sources/stacks-algebraic-stacks-morphisms/pinned-original-source.tex, source lines 8982--8982. Navigation pointer only: the counted text refers to the Stacks item morphisms-lemma-flat-base-change-scheme-theoretic-image, which belongs to another chapter of the project and is not part of this card. The pointer carries the target's own label as its number because this chapter's staged snapshot does not contain the other chapter's tag table; its content is not reproduced here.\par\medskip
\ifdefined\hypertarget\hypertarget{sourcepointer-4}{}\fi
\noindent\textit{Source reference.} Original reference morphisms-section-scheme-theoretic-image: sources/stacks-algebraic-stacks-morphisms/pinned-original-source.tex, source lines 8944--8944. Navigation pointer only: the counted text refers to the Stacks item morphisms-section-scheme-theoretic-image, which belongs to another chapter of the project and is not part of this card. The pointer carries the target's own label as its number because this chapter's staged snapshot does not contain the other chapter's tag table; its content is not reproduced here.\par\medskip
\ifdefined\hypertarget\hypertarget{sourcepointer-5}{}\fi
\noindent\textit{Source reference.} Original reference spaces-morphisms-lemma-flat-base-change-scheme-theoretic-image: sources/stacks-algebraic-stacks-morphisms/pinned-original-source.tex, source lines 8955--8955. Navigation pointer only: the counted text refers to the Stacks item spaces-morphisms-lemma-flat-base-change-scheme-theoretic-image, which belongs to another chapter of the project and is not part of this card. The pointer carries the target's own label as its number because this chapter's staged snapshot does not contain the other chapter's tag table; its content is not reproduced here.\par\medskip
\ifdefined\hypertarget\hypertarget{sourcepointer-6}{}\fi
\noindent\textit{Source reference.} Original reference spaces-morphisms-section-scheme-theoretic-image: sources/stacks-algebraic-stacks-morphisms/pinned-original-source.tex, source lines 8946--8946. Navigation pointer only: the counted text refers to the Stacks item spaces-morphisms-section-scheme-theoretic-image, which belongs to another chapter of the project and is not part of this card. The pointer carries the target's own label as its number because this chapter's staged snapshot does not contain the other chapter's tag table; its content is not reproduced here.\par\medskip
\ifdefined\hypertarget\hypertarget{sourcepointer-7}{}\fi
\noindent\textit{Source reference.} Original reference stacks-properties-lemma-substacks-presentation: sources/stacks-algebraic-stacks-morphisms/pinned-original-source.tex, source lines 4851--4851. Navigation pointer only: the counted text refers to the Stacks item stacks-properties-lemma-substacks-presentation, which belongs to another chapter of the project and is not part of this card. The pointer carries the target's own label as its number because this chapter's staged snapshot does not contain the other chapter's tag table; its content is not reproduced here.\par\medskip
\begin{lemma}
\label{lemma-existence-plus-flat-base-change}
Let $f : \mathcal{X} \to \mathcal{Y}$ be a quasi-compact
morphism of algebraic stacks. Then formation of the scheme theoretic image
commutes with flat base change.
\end{lemma}

\begin{proof}
Let $\mathcal{Y}' \to \mathcal{Y}$ be a flat morphism of algebraic stacks.
Choose a scheme $V$ and a surjective smooth morphism $V \to \mathcal{Y}$.
Choose a scheme $V'$ and a
surjective smooth morphism $V' \to \mathcal{Y}' \times_\mathcal{Y} V$.
We may and do assume that $V = \coprod_{i \in I} V_i$ is a disjoint
union of affine schemes and that
$V' = \coprod_{i \in I} \coprod_{j \in J_i} V_{i, j}$
is a disjoint union of affine schemes with each $V_{i, j}$ mapping into $V_i$.
Let
\begin{enumerate}
\item $\mathcal{Z} \subset \mathcal{Y}$ be the scheme theoretic image of $f$,
\item $\mathcal{Z}' \subset \mathcal{Y}'$ be the scheme theoretic image
of the base change of $f$ by $\mathcal{Y}' \to \mathcal{Y}$,
\item $Z \subset V$ be the scheme theoretic image
of the base change of $f$ by $V \to \mathcal{Y}$,
\item $Z' \subset V'$ be the scheme theoretic image
of the base change of $f$ by $V' \to \mathcal{Y}$.
\end{enumerate}
If we can show that
(a) $Z = V \times_\mathcal{Y} \mathcal{Z}$,
(b) $Z' = V' \times_{\mathcal{Y}'} \mathcal{Z}'$, and
(c) $Z' = V' \times_V Z$
then the lemma follows: the inclusion
$\mathcal{Z}' \to \mathcal{Z} \times_\mathcal{Y} \mathcal{Y}'$
(Lemma \ref{lemma-factor-factor})
has to be an isomorphism because after base change by the surjective
smooth morphism $V' \to \mathcal{Y}'$ it is.

\medskip\noindent
Proof of (a). Set $R = V \times_\mathcal{Y} V$.
By Properties of Stacks, Lemma
\ref{stacks-properties-lemma-substacks-presentation}
the rule $\mathcal{Z} \mapsto \mathcal{Z} \times_\mathcal{Y} V$
defines a $1$-to-$1$ correspondence between closed substacks
of $\mathcal{Y}$ and $R$-invariant closed subspaces of $V$.
Moreover, $f : \mathcal{X} \to \mathcal{Y}$ factors through $\mathcal{Z}$
if and only if the base change
$g : \mathcal{X} \times_\mathcal{Y} V \to V$ factors through
$\mathcal{Z} \times_\mathcal{Y} V$.
We claim: the scheme theoretic image $Z \subset V$ of $g$
is $R$-invariant. The claim implies (a) by what we just said.

\medskip\noindent
For each $i$ the morphism $\mathcal{X} \times_\mathcal{Y} V_i \to V_i$
is quasi-compact and hence $\mathcal{X} \times_\mathcal{Y} V_i$ is
quasi-compact. Thus we can choose an affine scheme $W_i$ and a
surjective smooth morphism $W_i \to \mathcal{X} \times_\mathcal{Y} V_i$.
Observe that $W = \coprod W_i$ is a scheme endowed with
a smooth and surjective morphism $W \to \mathcal{X} \times_\mathcal{Y} V$
such that the composition $W \to V$ with $g$ is quasi-compact.
Let $Z \to V$ be the scheme theoretic image of $W \to V$, see
Morphisms, Section
\ref{morphisms-section-scheme-theoretic-image} and
Morphisms of Spaces, Section
\ref{spaces-morphisms-section-scheme-theoretic-image}.
It follows from Lemma \ref{lemma-cover-upstairs}
that $Z \subset V$ is the scheme theoretic image of $g$.
To show that $Z$ is $R$-invariant we claim that both
$$
\text{pr}_0^{-1}(Z), \text{pr}_1^{-1}(Z) \subset R = V \times_\mathcal{Y} V
$$
are the scheme theoretic image of $\mathcal{X} \times_\mathcal{Y} R \to R$.
Namely, we first use Morphisms of Spaces, Lemma
\ref{spaces-morphisms-lemma-flat-base-change-scheme-theoretic-image}
to see that $\text{pr}_0^{-1}(Z)$ is the scheme theoretic image
of the composition
$$
W \times_{V, \text{pr}_0} R = W \times_\mathcal{Y} V \to
\mathcal{X} \times_\mathcal{Y} R \to R
$$
Since the first arrow here is surjective and smooth we see that
$\text{pr}_0^{-1}(Z)$ is the scheme theoretic image of
$\mathcal{X} \times_\mathcal{Y} R \to R$. The same argument applies
that $\text{pr}_1^{-1}(Z)$. Hence $Z$ is $R$-invariant.

\medskip\noindent
Statement (b) is proved in exactly the same way as one proves (a).

\medskip\noindent
Proof of (c). Let $Z_i \subset V_i$ be the scheme theoretic image
of $\mathcal{X} \times_\mathcal{Y} V_i \to V_i$ and let
$Z_{i, j} \subset V_{i, j}$ be the scheme theoretic image of
$\mathcal{X} \times_\mathcal{Y} V_{i, j} \to V_{i, j}$.
Clearly it suffices to show that the inverse image of $Z_i$
in $V_{i, j}$ is $Z_{i, j}$. Above we've seen that
$Z_i$ is the scheme theoretic image of $W_i \to V_i$
and by the same token $Z_{i, j}$ is the scheme theoretic
image of $W_i \times_{V_i} V_{i, j} \to V_{i, j}$.
Hence the equality follows from the case of schemes
(Morphisms, Lemma
\ref{morphisms-lemma-flat-base-change-scheme-theoretic-image})
and the fact that $V_{i, j} \to V_i$ is flat.
\end{proof}

\end{document}
